\section{Discussion}

% Claims C03-C10
This controlled audit found no complete transfer of the proposed BuNN
anti-collapse advantage to the specified ABIDE-I functional-connectome
pipeline. Learned bundle transport did not improve the primary performance-
density curve relative to GCN, did not exceed the strongest regularized
connectome baseline, and did not show the proposed matched-anchor effective-
rank advantage. The result is scientifically useful because the design
isolated propagation from backbone, input, tuning-budget, and capacity
differences.

\subsection{Why transport may not help in this setting}

Two properties distinguish this task from the settings that motivate bundle
transport. First, functional-connectome graphs are dense even after thresholding.
Second, the feature attached to each node is its full connectivity row, so each
node already contains global information before propagation. Diffusion may
therefore average globally informative profiles rather than deliver missing
long-range context. The downward descriptive density curves are consistent
with density-dependent aggregation degradation, but do not alone prove a
general over-smoothing mechanism.

\subsection{Importance of the learned-local control}

BuNN contains additional learned map-generator parameters. A comparison with
GCN alone could confuse transport with added node-wise capacity. The
learned-local control retained this capacity while removing inter-node
diffusion. Its high raw effective rank showed why unaligned or unmatched rank
curves could falsely appear to support bundle transport. The matched-anchor
analysis instead tested what changed when inter-node transport was added.

\subsection{Heterogeneity and uncertainty}

The primary BuNN-GCN interval crossed zero, and its point estimate changed sign
under participant weighting. Site exclusions and final seeds also changed the
magnitude and sometimes the direction of the small contrast. These analyses do
not reveal a hidden positive result; they show that a small average effect is
unstable relative to multi-site and optimization variation. In contrast, the
negative BuNN-versus-elastic-net result and matched-anchor effective-rank
finding were stable across all leave-one-site exclusions.

\subsection{Limitations}

The conclusions are conditional on ABIDE-I, C-PAC filtered data without global-
signal regression, the AAL-116 atlas, Pearson/Fisher-$z$ connectomes, positive-
edge thresholding, connectivity-profile node features, the fixed density grid,
the shared backbone, and the equal tuning budget. Functional correlations do
not identify anatomy or causal information flow. The study has no locked
external validation cohort, and site sample sizes vary substantially. Runtime
comparisons depend on this implementation and hardware. Finally, the
representation diagnostics measure computational behavior; they do not show
that a representation change caused a predictive change.

\subsection{Future work}

The next scientifically clean extension would freeze the complete ABIDE-I code
and evaluate it on a preprocessing-compatible ABIDE-II cohort. Separate
pre-registered sensitivity studies could test global-signal regression,
negative-edge handling, another atlas, or node features that do not already
encode the full connectome. Such studies should remain distinct from the
present confirmatory result and should not search across datasets for a
favorable diagnosis.
